\subsection{Extremal points and barycenters}

PROPOSITION 3. --- Let K be a compact convex subset of a Hausdorff locally convex space E, x a point of K. Every measure \(\mu\) on K, positive, of total mass 1, and admitting x as barycenter, is in the vague closure of the set of discrete positive measures of total mass 1 that admit x as barycenter.

Let U be a neighborhood of \(\mu\) for the vague topology; we can suppose that U consists of the measures \(\nu\) on K such that

\[
\left| \mu (f _ {i}) - \nu (f _ {i}) \right| \leqslant \delta\tag{1}
\]

for a finite number of functions \(f_{i} \in \mathcal{C}(\mathrm{K}; \mathbf{C})\) ( \(1 \leqslant i \leqslant p\) ) and a number \(\delta > 0\) . For every point \(a \in K\) , there exists a closed convex neighborhood \(V_{a}\) of 0 in E such that

\[
\left| f _ {i} (\mathbf {y}) - f _ {i} (\mathbf {a}) \right| \leqslant \delta / 2\tag{2}
\]

for \(1 \leqslant i \leqslant p\) and for every \(\mathbf{y} \in \mathrm{W}_{\mathbf{a}} = \mathrm{K} \cap (\mathbf{a} + \mathrm{V}_{\mathbf{a}})\) . Since \(\mathbf{K}\) is compact, there exists a finite number of points \(\mathbf{a}_j\) ( \(1 \leqslant j \leqslant r\) ) of \(\mathbf{K}\) such that the \(\mathrm{W}_{\mathbf{a}_j}\) form a covering of \(\mathbf{K}\) ( \(1 \leqslant j \leqslant r\) ). Consider a continuous partition of unity \((g_j)_{1 \leqslant j \leqslant r}\) on \(\mathbf{K}\) , subordinate to the covering \((\mathrm{W}_{\mathbf{a}_j})\) , and set \(\alpha_j = \mu(g_j)\) for all \(j\) ; if \(\alpha_j \neq 0\) set \(\mu_j = \alpha_j^{-1} g_j \cdot \mu\) , and if \(\alpha_j = 0\) set \(\mu_j = \varepsilon_{\mathbf{a}_j}\) . Each of the measures \(\mu_j\) is positive, of total mass 1, and its support is contained in the compact convex set \(\mathrm{W}_{\mathbf{a}_j}\) ; moreover, by definition,

\[
\mu = \sum_ {j = 1} ^ {r} \alpha_ {j} \mu_ {j}\tag{3}
\]

since \(g_{j} \cdot \mu = 0\) if \(\mu(g_{j}) = 0\) ; the \(\alpha_{j}\) are \(\geqslant 0\) and

\[
\sum_ {j = 1} ^ {r} \alpha_ {j} = \sum_ {j = 1} ^ {r} \mu (g _ {j}) = \mu \left(\sum_ {j = 1} ^ {r} g _ {j}\right) = \mu (1) = 1.
\]

Let \(\mathbf{x}_j\) be the barycenter of \(\mu_j\) , which belongs to \(\mathrm{W}_{\mathbf{a}_j}\) (No.~1, Prop. 1), and consider the discrete measure \(\nu = \sum_{j=1}^{r} \alpha_j \varepsilon_{\mathbf{x}_j}\) ; it is positive and of total mass 1, and its barycenter is \(\sum_{j=1}^{r} \alpha_j \mathbf{x}_j\) , which is also the barycenter of \(\mu\) by virtue of (3), thus is equal to \(\mathbf{x}\) . Moreover, by (2), \(|f_i(\mathbf{y}) - f_i(\mathbf{a}_j)| \leqslant \delta/2\) for all \(\mathbf{y} \in \mathrm{W}_{\mathbf{a}_j}\) and for all \(i\) , whence, since \(\operatorname{Supp}(\mu_j) \subset \mathrm{W}_{\mathbf{a}_j}\) , \(|\mu_j(f_i) - f_i(\mathbf{a}_j)| \leqslant \delta/2\) for \(1 \leqslant i \leqslant p\) . On the other hand, since \(\mathbf{x}_j \in \mathrm{W}_{\mathbf{a}_j}\) , one also has

\[
\left| \varepsilon_ {\mathbf {x} _ {j}} (f _ {i}) - f _ {i} (\mathbf {a} _ {j}) \right| \leqslant \delta / 2
\]

for \(1 \leqslant i \leqslant p\) , whence \(|\mu_j(f_i) - \varepsilon_{\mathbf{x}_j}(f_i)| \leqslant \delta\) for all \(i\) and \(j\) . Since the \(\alpha_j\) are \(\geqslant 0\) and have sum 1, it follows from (3) and the definition of \(\nu\) that \(\nu\) satisfies the inequality (1).

Q.E.D.

COROLLARY. --- Let \(K'\) be a compact subset of K such that K is the closed convex envelope of \(K'\) . In order that \(x \in K'\) be an extremal point of K, it is necessary and sufficient that \(\varepsilon_{x}\) be the only positive measure on \(K'\) , of total mass 1, having x as barycenter.

Suppose \(\mathbf{x}\) is an extremal point of \(K\) ; to prove that \(\varepsilon_{\mathbf{x}}\) is the only positive measure on \(K'\) , of total mass 1, having \(\mathbf{x}\) as barycenter, it suffices, by Prop. 3, to see that the set of discrete measures \(\nu\) on \(K'\) that are positive, of total mass 1, and have \(\mathbf{x}\) as barycenter, reduces to \(\varepsilon_{\mathbf{x}}\) . But such a measure \(\nu\) is of the form \(\sum_{i=1}^{r} \lambda_i \varepsilon_{\mathbf{x}_i}\) with \(\lambda_i > 0\) for \(1 \leqslant i \leqslant r\) and \(\sum_{i=1}^{r} \lambda_i = 1\) , and the hypothesis that \(\mathbf{x}\) is the barycenter of \(\nu\) may be written \(\mathbf{x} = \sum_{i=1}^{r} \lambda_i \mathbf{x}_i\) . Since \(\mathbf{x}\) is extremal, this implies that \(\mathbf{x}_i = \mathbf{x}\) for all \(i\) , whence \(\nu = \varepsilon_{\mathbf{x}}\) .

Conversely, let us assume that \(\varepsilon_{x}\) is the only positive measure on \(K'\) , of total mass 1, having x as barycenter, and let us show that x is extremal. In the contrary case, there would exist two distinct points \(x', x''\) of K and a real number \(\lambda\) such that \(0 < \lambda < 1\) and \(\mathbf{x} = \lambda\mathbf{x}' + (1 - \lambda)\mathbf{x}''\) . By Prop. 1, \(x'\) (resp. \(x''\) ) is the barycenter of a positive measure \(\mu'\) (resp. \(\mu''\) ) on \(K'\) of total mass 1. Then x is the barycenter of \(\lambda\mu' + (1 - \lambda)\mu''\) . Therefore \(\lambda\mu' + (1 - \lambda)\mu'' = \varepsilon_{x}\) . Therefore \(\mu'\) and \(\mu''\) are proportional to \(\varepsilon_{x}\) , whence \(x' = x'' = x\) , which is absurd.

THEOREM 1 (Choquet). --- Let E be a Hausdorff locally convex space over R, K a metrizable compact convex subset of E, and M the set of extremal points of K. The set M is the intersection of a countable family of open sets in K, and every point of K is the barycenter of a measure \(\mu\) on K such that \(\mu(K - M) = 0\) .

To prove the first assertion, denote by I the interval {[}0,1{]} of \(\mathbf{R}\) ; since K is compact and metrizable, so is \(\mathrm{K} \times \mathrm{K} \times \mathrm{I}\) ; the subset U of \(\mathrm{K} \times \mathrm{K} \times \mathrm{I}\) formed by the triples \((\mathbf{x}, \mathbf{y}, \lambda)\) such that \(\mathbf{x} \neq \mathbf{y}\) and \(0 < \lambda < 1\) is open in \(\mathrm{K} \times \mathrm{K} \times \mathrm{I}\) , therefore there exists a sequence \((\mathrm{F}_n)\) of closed sets in \(\mathrm{K} \times \mathrm{K} \times \mathrm{I}\) whose union is U (GT, IX, §2, No.~5, Prop. 7). The mapping \(q: \mathrm{K} \times \mathrm{K} \times \mathrm{I} \to \mathrm{K}\) defined by \(q(\mathbf{x}, \mathbf{y}, \lambda) = \lambda \mathbf{x} + (1 - \lambda) \mathbf{y}\) is continuous and, by the definition of extremal point, \(\mathrm{K} - \mathrm{M} = q(\mathrm{U}) = \bigcup_{n} q(\mathrm{F}_n)\) ; but \(\mathrm{F}_n\) is compact since it is closed in \(\mathrm{K} \times \mathrm{K} \times \mathrm{I}\) , therefore \(q(\mathrm{F}_n)\) is compact and therefore closed in K; the set \(\mathrm{U}_n = \mathrm{K} - q(\mathrm{F}_n)\) is therefore open in K, and we have \(\mathrm{M} = \bigcap \mathrm{U}_n\) .

In the continuation of the proof, we shall denote by u a continuous and convex numerical function on K, by \(G \subset E \times R\) the graph of u, and by S the closed convex envelope of G in \(E \times R\) .

Lemma 1. --- Let \(\overline{u}\) be the lower envelope of the continuous affine linear functions on E that are \(\geqslant u\) on K. Then S is the set of points \((\mathbf{a}, b) \in \mathrm{E} \times \mathbf{R}\) such that \(\mathbf{a} \in \mathrm{K}\) and \(u(\mathbf{a}) \leqslant b \leqslant \overline{u}(\mathbf{a})\) .

By the Hahn-Banach theorem, for \((\mathbf{a}, b)\) to belong to S, it is necessary and sufficient that \(h(\mathbf{a}, b) \geqslant 0\) for every continuous affine linear function h on \(E \times R\) such that \(h(\mathbf{x}, u(\mathbf{x})) \geqslant 0\) for \(x \in K\) . Setting \(h(\mathbf{x}, t) = f(\mathbf{x}) - \lambda t\) , where f is a continuous affine linear function on E, we see that the relation \((\mathbf{a}, b) \in S\) is equivalent to the following property: the relation

\[
f (\mathbf {x}) \geqslant \lambda u (\mathbf {x}) \quad \text { for   all } \mathbf {x} \in K\tag{4}
\]

implies

\[
f (\mathbf {a}) \geqslant \lambda b.\tag{5}
\]

First set \(\lambda = 0\) ; the fact that (4) implies (5) for every continuous affine linear function \(f\) on \(\mathbf{E}\) is equivalent to the relation \(\mathbf{a} \in \mathbf{K}\) by the Hahn-Banach theorem. Next, set \(\lambda = -1\) and replace \(f\) by \(-f\) ; then the relation \(f(\mathbf{x}) \leqslant u(\mathbf{x})\) in \(\mathbf{K}\) must imply \(f(\mathbf{a}) \leqslant b\) . But since \(u\) is convex and continuous on \(\mathbf{K}\) , \(u(\mathbf{a})\) is equal to the supremum of the \(f(\mathbf{a})\) for the continuous affine linear functions \(f\) on \(\mathbf{E}\) such that \(f(\mathbf{x}) \leqslant u(\mathbf{x})\) on \(\mathbf{K}\) (TVS, II, §5, No.~4, Prop. 5); we thus obtain the relation \(b \geqslant u(\mathbf{a})\) . Finally, set \(\lambda = 1\) ; to say that (4) implies (5) then means, by definition, that \(b \leqslant \overline{u}(\mathbf{a})\) . This proves the lemma, since the relation (4) (resp. (5)) is equivalent to the one obtained by multiplying both sides by a scalar \(>0\) .

Lemma 2. --- If u is strictly convex on K, then \(u(\mathbf{x}) < \overline{u}(\mathbf{x})\) for every non-extremal point x of K.

For, there then exist two distinct points y,z of K such that \(\mathbf{x} = (\mathbf{y} + \mathbf{z})/2\) , whence \(u(\mathbf{x}) < (u(\mathbf{y}) + u(\mathbf{z}))/2\) since u is strictly convex. If f is an affine linear function on E, then \(f(\mathbf{x}) = (f(\mathbf{y}) + f(\mathbf{z}))/2\) ; applying this relation to the continuous affine linear functions f that are \(\geqslant u\) on K, we obtain

\[
\overline {{{u}}} (\mathbf {x}) \geqslant \left(\overline {{{u}}} (\mathbf {y}) + \overline {{{u}}} (\mathbf {z})\right) / 2 \geqslant \left(u (\mathbf {y}) + u (\mathbf {z})\right) / 2 > u (\mathbf {x}).
\]

These lemmas established, we therefore (Lemma 1) have \((\mathbf{a},\overline{u} (\mathbf{a}))\in S\) for all \(\mathbf{a}\in K\) . Since G is compact, being the image of K under the continuous mapping \(\mathbf{x}\mapsto (\mathbf{x},u(\mathbf{x}))\) , there exists, by Prop. 1 of No.~1, a positive measure \(\nu\) on G, of total mass 1, having \((\mathbf{a},\overline{u} (\mathbf{a}))\) as barycenter. Since the restriction of the projection \(\mathrm{pr}_1\) to G is a homeomorphism of G onto K, the measure \(\nu\) may be transported by means of this homeomorphism, which yields a measure \(\mu\) on K (positive and of total mass 1) such that

\[
\mathbf {a} = \int \mathbf {x} d \mu (\mathbf {x}) \quad \mathrm{and} \quad \overline {{{u}}} (\mathbf {a}) = \int u (\mathbf {x}) d \mu (\mathbf {x}).\tag{6}
\]

The first of the relations (6) means that a is the barycenter of \(\mu\) . The function \(\overline{u}\) is upper semi-continuous and bounded on K, hence is \(\mu\) -integrable ( \(\S4\) , No.~4, Cor. 1 of Prop. 5); moreover, since the function \(-\overline{u}\) is by definition convex, we have, by the Cor. of Prop. 2 of No.~1,

\[
\overline {{{u}}} (\mathbf {a}) \geqslant \int \overline {{{u}}} (\mathbf {x}) d \mu (\mathbf {x}),\tag{7}
\]

whence, on comparing with the second of the formulas (6),

\[
\int u (\mathbf {x}) d \mu (\mathbf {x}) \geqslant \int \overline {{{u}}} (\mathbf {x}) d \mu (\mathbf {x}).\tag{8}
\]

But since \(u(\mathbf{x}) \leqslant \overline{u}(\mathbf{x})\) on K, the relation (8) implies that \(u(\mathbf{x}) = \overline{u}(\mathbf{x})\) almost everywhere for \(\mu\) . Taking into account Lemma 2, one sees that Theorem 1 will be proved once the following lemma has been established:

Lemma 3. --- Let E be a Hausdorff locally convex space over R, and K a metrizable compact convex subset of E. Then, there exists a strictly convex numerical function on K.

For, the Banach space \(\mathcal{C}(\mathrm{K};\mathbf{R})\) is separable (GT, X, §3, No.~3, Th. 1), therefore so is the subspace \(\mathcal{A}\) of \(\mathcal{C}(\mathrm{K};\mathbf{R})\) formed by the restrictions to K of the continuous affine linear functions on E. Thus let \((h_n)\) be a sequence dense in \(\mathcal{A}\) , and let \(\alpha_{n} > \sup_{\mathbf{x}\in \mathbb{K}}|h_n(\mathbf{x})|\) . Then each of the functions \(h_n^2 / n^2\alpha_n^2\) is convex in K (TVS, II, §2, No.~8, Examples), and the series with general term \(h_n^2 / n^2\alpha_n^2\) is normally convergent, therefore its sum \(u\) is continuous and convex in K. It remains to see that \(u\) is strictly convex, and for this it suffices to prove that for any two distinct points \(\mathbf{x},\mathbf{x}'\) of K, there is an integer \(n\) such that the restriction of \(h_n^2\) to the segment with endpoints \(\mathbf{x},\mathbf{x}'\) is strictly convex; but for this it suffices that \(h_n(\mathbf{x})\neq h_n(\mathbf{x}')\) (loc. cit.). Now, there exists a function \(h\in \mathcal{A}\) such that \(h(\mathbf{x})\neq h(\mathbf{x}')\) (TVS, II, §4, No.~1, Cor. 1 of Prop. 2) and since the sequence \((h_n)\) is dense in \(\mathcal{A}\) , there exists an \(n\) such that \(h_n(\mathbf{x})\neq h_n(\mathbf{x}')\) .

Q.E.D.

COROLLARY. --- Let E be a Hausdorff locally convex space over R, C a proper convex cone in E with vertex 0, that is complete and metrizable for the uniform structure induced by the weakened uniform structure of E. Let M be the union of the extremal generators of C. For every \(x \in C\) there exist a compact convex subset K of C and a measure \(\lambda \geqslant 0\) on K of total mass 1, such that \(K - (M \cap K)\) is \(\lambda\) -negligible and the barycenter of \(\lambda\) is equal to x.

For, x belongs to a cap K of C (TVS, II, §7, No.~2, Prop. 5), and M∩K contains the set of extremal points of K (loc. cit., Cor. 1 of Prop. 4). It then suffices to apply Th. 1.

